\section{Arquitetura funcional e alimentação}
\label{sec:arquitetura}

A Figura~\ref{fig:arquitetura} diferencia as rotinas térmicas utilizadas na
classificação da contagem de transições do SW-420. Os blocos internos
representam funções executadas na mesma ESP32, e não microcontroladores
distintos. O treinamento ocorre previamente no computador; o dispositivo
executa somente a inferência do modelo exportado.

\begin{figure}[htbp]
\centering
\caption{Arquitetura implementada: aquisição, inferência local e serviços externos}
\label{fig:arquitetura}
\resizebox{\textwidth}{!}{\begin{tikzpicture}[x=1cm,y=1cm]
\node[block,text width=3.9cm] (dht) at (-4.8,0.9) {DHT22\\Temperatura e umidade do ar};
\node[block,text width=3.9cm] (probe) at (0,0.9) {Termopar K + MAX6675\\Temperatura no ponto da sonda};
\node[block,text width=3.9cm] (sw) at (4.8,0.9) {SW-420\\Saída digital DO};
\node[block,text width=7.1cm] (thermal) at (-2.4,-1.9) {Aquisição térmica e validação\\DHT22 e sonda independentes};
\node[block,text width=3.9cm] (count) at (4.8,-1.9) {Interrupção CHANGE\\Contagem por janela e DO};
\node[block,text width=7.1cm] (ai) at (-2.4,-3.7) {Preparação de quatro atributos\\Random Forest em C++\\Vibração não integra o modelo};
\node[block,text width=3.9cm] (serial) at (4.8,-3.7) {Saída serial\\Medições e classificação\\115200 baud};
\node[block,text width=10.5cm] (packet) at (-0.4,-5.6) {JSON: medições + diagnóstico embarcado\\Comunicação Wi-Fi / HTTP POST quando conectado};
\begin{scope}[on background layer]
\node[draw=black!50,dashed,rounded corners,fill=blue!1,inner sep=0.35cm,
fit=(thermal)(count)(ai)(serial)(packet)] (esp) {};
\end{scope}
\node[font=\footnotesize,fill=white] at ($(esp.north)+(1.7,0)$) {Rotinas na ESP32 Heltec V2};
\node[block,text width=7.2cm] (api) at (-1.5,-7.7) {API FastAPI no computador\\Validação e consulta dos registros};
\node[block,text width=3.1cm] (csv) at (5.0,-7.7) {Arquivo CSV\\Persistência local};
\node[block,text width=8cm] (web) at (-1.5,-9.7) {Interface web React\\Medições, classificação, frescor e histórico};
\draw[flow] (dht.south) -- ++(0,-0.45) -| ([xshift=-1.6cm]thermal.north);
\draw[flow] (probe.south) -- ++(0,-0.45) -| ([xshift=1.6cm]thermal.north);
\draw[flow] (sw) -- (count);
\draw[flow] (thermal) -- (ai);
\draw[flow] (ai.east) -- (serial.west);
\draw[flow] (count) -- (serial);
\draw[flow] (ai.south) -- ++(0,-0.3) -| ([xshift=-1.5cm]packet.north);
\draw[flow] (count.east) -- ++(0.65,0) |- (packet.east);
\draw[flow] (packet.south) -- (api.north) node[midway,left,edge label] {POST};
\draw[flow] (api) -- (csv);
\draw[{Latex[length=2.2mm]}-{Latex[length=2.2mm]},line width=0.65pt]
  (api.south) -- (web.north) node[midway,right,edge label] {GET / resposta JSON};
\end{tikzpicture}
}
\source{elaboração própria a partir do firmware, da API e do frontend do MS-Monitor}
\end{figure}

\begin{table}[htbp]
\centering
\caption{Responsabilidade dos módulos e limites da implementação}
\label{tab:modulos}
\small
\begin{tabularx}{\textwidth}{L{3cm}Y}
\toprule
Módulo & Função no protótipo\\
\midrule
Aquisição & Ler DHT22 e MAX6675; capturar mudanças do sinal digital de vibração por interrupção.\\
Pré-processamento & Validar as entradas, preservar indicadores de falha e calcular a diferença entre temperaturas.\\
IA embarcada & Executar o Random Forest exportado; produzir classe, índice, confiança e duração da chamada.\\
Interface local & Registrar medições e resultado na serial a 115200 baud, inclusive sem Wi-Fi.\\
Comunicação & Realizar uma tentativa de POST por leitura quando houver Wi-Fi; tratar o código de retorno.\\
API e armazenamento & Validar o contrato, registrar cada POST aceito em CSV e disponibilizar consultas.\\
Visualização & Exibir medições, origem da classificação, frescor dos dados e histórico no navegador.\\
Atuação & Não implementada; não há comando de ventilador no firmware atual.\\
\bottomrule
\end{tabularx}
\source{elaboração própria a partir do código do projeto}
\end{table}

\subsection{Caminhos de energia e de retorno}

A alimentação adotada é USB, com entrada nominal de 5~V na placa e
distribuição de 3,3~V aos módulos sensores. O terra é comum, conforme
a Figura~\ref{fig:alimentacao}. As linhas de alimentação e GND não representam
o fluxo lógico de dados. O termopar é passivo e se conecta às entradas
próprias do MAX6675; não recebe alimentação de 3,3~V na ponta sensora.

\begin{figure}[htbp]
\centering
\caption{Distribuição de alimentação adotada e referência de terra comum}
\label{fig:alimentacao}
\resizebox{\textwidth}{!}{\begin{tikzpicture}[x=1cm,y=1cm]
\node[block,text width=4.5cm] (usb) at (0,0) {Fonte / porta USB\\Entrada nominal de 5 V};
\node[block,text width=5cm] (board) at (0,-1.8) {Heltec V2\\Entrada USB e regulação da placa};
\node[block,text width=3.8cm] (dht) at (-5,-4.6) {DHT22\\VCC = 3,3 V};
\node[block,text width=3.8cm] (sw) at (0,-4.6) {Módulo SW-420\\VCC = 3,3 V};
\node[block,text width=3.8cm] (max) at (5,-4.6) {MAX6675\\VCC = 3,3 V};
\draw[flow] (usb) -- (board);
\draw[line width=1pt] (-5,-3.2) -- (5,-3.2);
\draw[flow] (board.south) -- (0,-3.2);
\node[font=\small,fill=white,anchor=south] at (-2.6,-3.2) {Barramento de 3,3 V};
\foreach \target/\xx in {dht/-5,sw/0,max/5}{
  \draw[flow] (\xx,-3.2) -- (\target.north);
  \draw[line width=0.7pt] (\target.south) -- (\xx,-6);
}
\draw[line width=0.7pt] (-5.5,-6) -- (7.2,-6);
\draw[line width=0.7pt] (board.east) -- (7.2,-1.8) -- (7.2,-6);
\node[font=\small,fill=white,anchor=north] at (0,-6) {GND comum: referência elétrica e caminho de retorno};
\node[font=\footnotesize,align=center,text width=13cm] at (0,-7.1)
{O termopar é passivo: liga-se a T+ / T-- do MAX6675.\\
O desenho não representa fios de dados nem um circuito de acionamento.};
\end{tikzpicture}
}
\source{elaboração própria a partir das ligações documentadas; diagrama funcional, não esquema completo da placa}
\end{figure}

Os sinais que chegam aos GPIOs devem ser compatíveis com a lógica de
3,3~V da ESP32. A placa utilizada possui recursos adicionais de comunicação
e visualização, mas LoRa, Bluetooth e OLED não são utilizados pelo firmware
analisado \cite{heltecV2,espressif2026}. O GPIO~21, associado ao controle
de Vext nessa montagem, não é destinado aos sensores.

\begin{table}[htbp]
\centering
\caption{Mapa de conexões utilizado no firmware}
\label{tab:pinos}
\begin{tabularx}{\textwidth}{L{4.0cm}L{2.3cm}Y}
\toprule
Sinal & GPIO & Interface e direção em relação à ESP32\\
\midrule
DHT22 --- DAT & 17 & Protocolo digital temporizado, bidirecional.\\
SW-420 --- DO & 22 & Entrada digital; interrupção em \texttt{CHANGE}.\\
MAX6675 --- SCK & 13 & Saída de relógio gerada por software.\\
MAX6675 --- CS & 23 & Saída de seleção do conversor.\\
MAX6675 --- SO & 39 & Entrada de dados; GPIO de entrada.\\
VCC dos módulos & --- & Barramento de 3,3~V.\\
GND dos módulos & --- & Retorno comum ao GND da placa.\\
\bottomrule
\end{tabularx}
\source{elaboração própria a partir das constantes e rotinas do firmware}
\end{table}

\subsection{Grandezas elétricas e dimensionamento pendente}

O dimensionamento final requer a identificação da fonte USB, a corrente
dos módulos efetivamente montados e os picos durante a comunicação. Esses
valores não foram determinados pelos testes de software. Assim, não se
apresenta uma soma de correntes hipotéticas como consumo medido do protótipo.
A potência de entrada deve ser obtida no mesmo ponto de medição:
\begin{equation}
P_{\mathrm{USB}} = V_{\mathrm{USB}} I_{\mathrm{USB}}.
\end{equation}
Já a potência entregue apenas aos sensores no barramento de 3,3~V é
\begin{equation}
P_{\mathrm{sensores}} =
3{,}3\left(I_{\mathrm{DHT22}}+I_{\mathrm{MAX6675}}+I_{\mathrm{SW420}}\right),
\end{equation}
com correntes em ampères e potência em watts. Essa expressão não inclui
todo o consumo da placa, nem as perdas de regulação.

Como exemplo exclusivamente didático, para um resistor ideal de
$10~\mathrm{k}\Omega$ submetido a 3,3~V, a Lei de Ohm fornece
\begin{equation}
I=\frac{V}{R}=\frac{3{,}3}{10\,000}
=0{,}00033~\mathrm{A}=0{,}33~\mathrm{mA}.
\end{equation}
Esse resistor não é apresentado como componente identificado da montagem.
O exemplo apenas ilustra a relação entre tensão, resistência e corrente.
O plano de testes da Seção~\ref{sec:avaliacao} prevê a medição da entrada
USB e do barramento de 3,3~V antes de concluir sobre a suficiência da fonte.
\FloatBarrier
